\section{Security Model}
\label{sec:security-model}

This section instantiates the security claims for \TW{}. The root tag is first
treated as the output of the wrapper $\HTW$: adversaries supply full inputs or
candidate openings, and the games ask whether the tag collides or can be hit as
a preimage. The same construction is also treated as a nonce-based AEAD: the
\textsf{INT-RUP} game gives integrity even when a plaintext-computing interface
is exposed, while all-in-one AE and mu-ind record the baseline
confidentiality and authenticity notions.

\Cref{tab:security-notions} summarizes the proof obligations before the formal
games. The loss terms named there are defined in \Cref{sec:loss-terms}, and the
implemented-parameter instantiations appear in \Cref{sec:concrete-security}.

\paragraph{Proof worlds.}
All final statements are public-permutation claims for \TW{} with \Keccakp{}
modeled as a uniformly random permutation. The proofs first establish random
oracle bounds for the flat duplex interface of \Cref{sec:ro-model}; the flat
sponge lift of \Cref{app:flat-sponge-lift} then transfers those bounds to the
public-permutation games of \Cref{sec:games-pp,sec:mu-setting}.

\paragraph{Scope of the security claims.}
In the AE-style distinguishing games, encryption queries are nonce-respecting
and the verification oracle returns only acceptance or rejection. The
\textsf{INT-RUP} game gives integrity under a plaintext-computing oracle, but it
does not give privacy under release of unverified plaintext or
plaintext-awareness security. None of the AE-style games gives nonce-misuse
resistance. The hash and committing games are adversarial-initialization games
over submitted inputs $(K, U, A, M)$ or ciphertext openings and impose no nonce
uniqueness condition. Across all games, the stated bounds are birthday-type,
single-stage, and leakage-free: beyond-birthday and side-channel security are
outside the model.
\Cref{sec:limitations} discusses the consequences of these exclusions and the
corresponding open problems, and \Cref{sec:permutation-cryptanalysis} reviews
the third-party cryptanalysis supporting the random-permutation modeling of
\Keccakp{}.

\begin{table}[t]
  \centering
  \footnotesize
  \setlength{\tabcolsep}{3pt}
  \begin{tabular}{@{}p{0.18\linewidth}p{0.34\linewidth}p{0.23\linewidth}p{0.17\linewidth}@{}}
    \toprule
    Notion & Adversary goal & Main result & Loss shape \\
    \midrule
    $\HTW$ collision
      & Find distinct full inputs $(K, U, A, M)$ with the same root tag.
      & \Cref{thm:coll}
      & Capacity lift, root collision, leaf collision \\
    $\HTW$ preimage
      & Find an input hitting a fixed target tag.
      & \Cref{thm:epre}
      & Capacity lift, root preimage \\
    $\HTW$ second preimage
      & Find a second preimage in the standard or everywhere
        \cite{RS04} sense.
      & \Cref{cor:second-preimage-rs04}
      & Inherited from collision resistance \\
    CMT / CMTD
      & Produce equal encryptions or one valid ciphertext opening under
        distinct committed inputs.
      & \Cref{cor:cmtds4,cor:committing-combined}
      & Capacity lift, root collision \\
    \textsf{INT-RUP}
      & Produce a fresh accepting ciphertext after encryption and
        plaintext-computing queries.
      & \Cref{thm:int-rup}
      & Capacity lift, key guess, leaf collision, tag guess \\
    AE / mu-ind
      & Distinguish real encryption and verification from random ciphertexts
        and replay-only verification, with nonce-respecting encryption queries.
      & \Cref{thm:mu-ind,cor:ae}
      & Capacity lift, key guess, leaf collision, tag guess \\
    \bottomrule
  \end{tabular}
  \caption{Proof obligations and the dominant loss shapes used for \TW{}.
    The table is only a roadmap; the formal games and resource accounting are
    defined in \Cref{sec:games-pp,sec:mu-setting,sec:resources}.}
  \label{tab:security-notions}
\end{table}

The proof dependency spine is as follows. \Cref{app:lemmas} establishes
transcript decoding, bridge injectivity, output-point separation, and the local
probabilistic accounting lemmas for key tests, leaf tag collisions, and root tag
guesses. \Cref{sec:root-tag-hash-security} uses those structural facts to build
a chronological sequence of final root candidates and apply a single root tag
hash bound; the hash and committing theorems differ only in which collision or
preimage event is reduced to that candidate sequence. \Cref{sec:aead-security}
uses the same facts for the random oracle mu-ind and \textsf{INT-RUP} bounds.
\Cref{app:flat-sponge-lift} transfers both branches to the public permutation
games defined below.

\subsection{Public Permutation Games}
\label{sec:games-pp}

\paragraph{Primitive convention.}
In every public permutation game below, the challenger samples
$\pi \sample \Perm(\bits^{1600})$ and exposes the direct primitive
interface $(\pi, \pi^{-1})$ to the adversary. The real-vs-ideal AE-style
games compare the real construction interface with the ideal interface
specified by the game over this shared primitive, following the
same-oracle convention of \cite{BT16}; the hash and committing games
below are one-world success probabilities over the same sampled
permutation. The simulator $(\Sim^{\ROflat},\Sim^{-1,\ROflat})$ of
\Cref{app:flat-sponge-lift} is a proof device, entering through the
lift from the intermediate random oracle model.

\paragraph{Validity convention.}
All construction interface inputs are checked against \TW{}'s public
syntax and domain before any construction computation or
construction-internal primitive lookup is performed. An encryption input
is valid when $U \in \bits^\nu$, $A$ is a byte string with
$\|A\| \leq A_{\max}$, and $M$ is a byte string with
$\|M\| \leq M_{\max}$. A verification or decryption input is valid when
$U \in \bits^\nu$, $A$ is a byte string with $\|A\| \leq A_{\max}$, and
the full ciphertext parses as $C \concat T$ with $\|C\| \leq M_{\max}$
and $T \in \bits^{\tauroot}$. In the hash and committing games, each
submitted input $(K, U, A, M)$ must be a valid encryption input including
$K \in \bits^k$, and a submitted CMTDs ciphertext must be a valid
verification/decryption ciphertext. Invalid encryption queries
are rejected and not recorded, invalid plaintext-computing queries return
$\bot$, invalid verification queries return $0$, and invalid committing-game
outputs make the game return $0$. Such rejected inputs contribute no
construction cost and create no construction transcripts. Throughout, an
\emph{accepted} encryption query is one that passes this validity check and any
game-specific nonce or user check.

\paragraph{Construction oracles.}
Whenever a game uses a secret key $K$ and permutation $\pi$, $\Enc_K[\pi]$
and $\Dec_K[\pi]$ are the algorithms of \Cref{alg:enc,alg:dec} with that
permutation. The verification oracle is
$\Ver_K[\pi](U, A, C, T) = [\Dec_K[\pi](U, A, C \concat T) \neq \bot]$.
For \textsf{INT-RUP}, $\PDec_K[\pi](U, A, C, T)$ is the plaintext recovered by
the decryption transcript before the final root-tag comparison; it returns
neither the recomputed tag nor the verification bit.

\paragraph{Hash games.}
The hash games for the wrapper $\HTW$ (\Cref{def:htw}) have no
challenger-sampled key. In the $s$-way collision game, all inputs are
adversary-supplied: the adversary outputs $s$ inputs, and the challenger
returns $1$ iff they are pairwise distinct, in the \TW{} domain, and share
a root tag:
\begin{align*}
  &\pi \sample \Perm(\bits^{1600}), \quad
   (K_i, U_i, A_i, M_i)_{i=1}^s \sample \mathcal{A}^{\pi,\, \pi^{-1}}, \\
  &\text{return $0$ unless the inputs are pairwise distinct and in the \TW{} domain}, \\
  &\text{return } \bigl[\HTW(K_1, U_1, A_1, M_1) = \cdots =
     \HTW(K_s, U_s, A_s, M_s)\bigr],
\end{align*}
with advantage $\AdvColl_\TW(\mathcal{A}) = \Pr[\text{the game returns }
1]$; no nonce uniqueness is imposed. The \cite{RS04} standard and
everywhere second-preimage notions are used on fixed input slices; in
\Cref{cor:second-preimage-rs04} they are reduced explicitly to this
two-way collision game.
For fixed byte lengths $a,\mu$, the standard preimage game samples
$X^* \sample \mathcal{X}_{a,\mu}$ independently of the primitive, computes
$T^* = \HTW(X^*)$, gives $T^*$ to
$\mathcal{A}^{\pi,\, \pi^{-1}}$, and returns $1$ iff $\mathcal{A}$ outputs a
valid input $X$ in the \TW{} domain with $\HTW(X) = T^*$; the advantage is
$\AdvPreSlice{a}{\mu}_\TW(\mathcal{A})$. The
everywhere preimage game fixes a target tag
$T^* \in \bits^{\tauroot}$ before the primitive sampling and returns $1$
iff $\mathcal{A}(T^*)$ outputs an input $X$ in the domain with
$\HTW(X) = T^*$, the advantage $\AdvePre_\TW$ taken as the worst case over
query-independent targets.

\paragraph{Committing games.}
The \cite{BH22} committing games are instantiated as in
\Cref{sec:commit-defs}, with no challenger-sampled key. The $s$-way
CMTDs-$\ell$ (decryption) game has the adversary output a ciphertext
$C \concat T$ and $s$ tuples $(K_i, U_i, A_i, M_i)$ with pairwise distinct
$\WiC_\ell$-images, and returns $1$ iff
$\Dec_{K_i}[\pi](U_i, A_i, C \concat T) = M_i$ for every $i$, with the eager
length rejection justified by \Cref{alg:dec}
($\Dec_K(U, A, C \concat T)$ returns $\bot$ or a message of length $\|C\|$);
its advantage is
$\Advcmtds{\ell}_\TW(\mathcal{A})$. The CMTs-$\ell$ (encryption) game has
the adversary output the $s$ tuples alone and returns $1$ iff the
encryptions $\Enc_{K_i}[\pi](U_i, A_i, M_i)$ are all equal; its advantage is
$\Advcmts{\ell}_\TW(\mathcal{A})$.

\paragraph{RUP integrity.}
The \textsf{INT-RUP} game samples $K \sample \bits^k$ and
$\pi \sample \Perm(\bits^{1600})$, exposes $\Enc_K[\pi]$,
$\PDec_K[\pi]$, $\Ver_K[\pi]$, and $(\pi, \pi^{-1})$, and maintains the replay
table $W$ of encryption outputs. Accepted encryption queries are
nonce-respecting: an encryption query is rejected if its nonce has appeared in a
previous accepted encryption query. Plaintext-computing and verification
queries may use arbitrary nonces. On an accepted encryption query $(U, A, M)$ the
game returns $\Enc_K[\pi](U, A, M)$, parsed and recorded as $(U, A, C, T)$ in $W$.
On a valid plaintext-computing query $(U, A, C, T)$ it returns
$\PDec_K[\pi](U, A, C, T)$, and on an invalid query it returns $\bot$. On a valid
verification query $(U, A, C, T)$ it returns $\Ver_K[\pi](U, A, C, T)$; if this bit
is $1$ and $(U, A, C, T) \notin W$ when the query is asked, the game records a
forgery. Invalid verification queries return $0$. The game returns $1$ iff a
forgery is recorded, and
\[
  \AdvintRUP_\TW(\mathcal{A}) =
  \Pr[\mathsf{G}^{\mathsf{INT\text{-}RUP}}_\TW(\mathcal{A}) \Rightarrow 1].
\]

\paragraph{All-in-one AE.}
The all-in-one authenticated encryption game~\cite{RS06} samples
$K \sample \bits^k$ and compares the real pair
$(\Enc_K[\pi], \Ver_K[\pi])$ with the ideal pair $(\Rand, \Replay)$.
On each accepted encryption query $(U, A, M)$, $\Rand$ samples a uniformly
random byte string of length $\|M\| + \tauroot/8$, parses it as
$C \concat T$ with $T \in \bits^{\tauroot}$, records $(U, A, C, T)$, and
returns $C \concat T$; $\Replay$ accepts exactly recorded tuples. Encryption
queries are nonce-respecting, and
\[
  \Advae_\TW(\mathcal{A}) =
  \bigl|
    \Pr[\mathcal{A}^{\Enc_K[\pi],\, \Ver_K[\pi],\, \pi,\, \pi^{-1}} \Rightarrow 1]
    - \Pr[\mathcal{A}^{\Rand,\, \Replay,\, \pi,\, \pi^{-1}} \Rightarrow 1]
  \bigr|.
\]
This all-in-one AE advantage is the single-user case of the multi-user
real-vs-ideal notion of \Cref{sec:mu-ind}; \Cref{cor:ae} derives it from
\Cref{thm:mu-ind}.

\subsection{Multi-User Setting}
\label{sec:mu-setting}

This subsection instantiates the mu-ind game of \cite{BT16}
(\Cref{sec:mu-ind}) for \TW{} in the public permutation model, under
the primitive convention of \Cref{sec:games-pp}. The challenger
samples a challenge bit $b \sample \bits$ and a permutation
$\pi \sample \Perm(\bits^{1600})$, and exposes $(\pi, \pi^{-1})$ as the
primitive interface in both worlds. The
adversary $\mathcal{A}$ creates users on demand via $\mathsf{New}$, and
the game maintains an active-user counter $v$. We write $D$ for an
execution-uniform cap on the number of $\mathsf{New}$ calls. Equivalently,
the game may sample independent keys $K_1,\dots,K_D$ at the start of the
experiment and let the $d$-th $\mathsf{New}$ call activate $K_d$; inactive
indices are never valid users.

The mu-ind construction oracles inherit the validity convention of
\Cref{sec:games-pp}. The encryption oracle on input $(d, U, A, M)$
rejects, without recording, if $(U, A, M)$ is outside the \TW{} encryption
domain, if $d \notin \{1, \dots, v\}$, or if $(d, U) \in V$; otherwise
it sets $C = \Enc_{K_d}[\pi](U, A, M)$ when $b = 1$, and samples a
uniformly random byte string $C$ of length $\|M\| + \tauroot/8$ when
$b = 0$. The game records
$(d, U)$ in $V$ and $(d, U, A, C)$ in $W$, and returns $C$. The
verification oracle on input $(d, U, A, C)$ rejects if $d \notin
\{1, \dots, v\}$ or if $(U, A, C)$ is not a valid \TW{}
verification/decryption input; if $(d, U, A, C) \in W$ it returns
$\mathsf{true}$; if $b = 0$ it returns $\mathsf{false}$; otherwise it
returns $[\Dec_{K_d}[\pi](U, A, C) \neq \bot]$. Nonce uniqueness is
enforced per user by the set $V$.

After querying these oracles, $\mathcal{A}$ outputs $b' \in \bits$.
The mu-ind advantage of $\mathcal{A}$ against \TW{} is
\[
  \Advmuind_\TW(\mathcal{A})
  = |2 \Pr[b' = b] - 1|.
\]
A valid verification query is \emph{non-replay} if its tuple is not in
$W$ when the query is asked; replay hits are answered by the recorded
transcript and are not counted in the verification-query resource below.

Resource caps split per potential user $d \leq D$: $N^{(d)}$,
$\qv^{(d)}$, $\nleaf^{(d)}$ are the per-user counts under
\Cref{sec:resources}, with inactive users contributing zero. The global
totals $N = \sum_{d=1}^D N^{(d)}$, $\qv = \sum_{d=1}^D \qv^{(d)}$,
$\nleaf = \sum_{d=1}^D \nleaf^{(d)}$ appear in the final bound: $N$
through the lift term $\Lift_c(\Ntot)$, and $\nleaf$ and $\qv$ through
the leaf-collision and verification-tag terms. $\Nprim$
remains a single shared count over the primitive interface.

The mu-ind treatment applies only to indistinguishability. The
committing security games already give the adversary control of all
$s$ keys, so the $s$-way CMTDs-$\ell$ statements of
\Cref{sec:games-pp} are themselves the multi-key committing statement
and need no separate $\mathsf{mu\text{-}cmt}$ notion.

\subsection{The TW128 Random Oracle Model}
\label{sec:ro-model}

The proofs of \Cref{sec:aead-security,sec:root-tag-hash-security} and
\Cref{cor:committing-combined} use an intermediate random oracle model. This
model keeps the same construction-level games but replaces every duplex output
by a lookup to a single random oracle through the \cite{BDPV11} bridge of
\Cref{sec:sponge}. The flat sponge lift of \Cref{app:flat-sponge-lift} then
converts the resulting bounds to the public-permutation form.

The model has a single random oracle $\ROflat : \bits^* \to
\bits^\infty$ over the unpadded $H$-input domain of the \cite{BDPV08}
simple-padded sponge interface, with lazy sampling: each distinct
input receives an independent infinite bit string, and repeated queries
return the same string. A direct adversary query is finite-output: on
input $(Y, \ell)$ with finite $\ell$, the oracle returns the first
$\ell$ bits of the lazy-sampled value for $Y$.

Construction calls do not query $\ROflat$ on arbitrary strings. Every root
or leaf duplex output is read from $\ROflat$ at the bridged flattened
transcript prefix for that duplex call, and the caller takes only the
requested output projection. Direct adversary queries to $\ROflat$ are the
only random oracle queries counted as direct oracle queries; construction
lookups are accounted for through the construction cost $N$ of
\Cref{sec:resources}.

For a well-formed \TW{} duplex transcript prefix $X$ in the language of
\Cref{def:wf-transcripts}, at one of the construction transcript lookup
points formalized by \Cref{lem:bridge-decode,tab:keyed-twout}, let
$\flatmap(X) = \bridge[1344, 1344](\mathsf{raw\_flat}(X))$ be the
bridged image of the native flattened input under the
\cite[Theorem~3]{BDPV11} preprocessing map of \Cref{sec:sponge}, and define
\[
  \RF(X) = \ROflat(\flatmap(X)).
\]
The construction interfaces $\Enc_K^{\RO}$, $\Dec_K^{\RO}$,
$\PDec_K^{\RO}$, and $\Ver_K^{\RO}$ are the \TW{} algorithms of
\Cref{sec:tw128} with every root or leaf duplex output replaced by the
corresponding requested projection of $\RF(X)$; the projection length is the
duplexing call's output parameter $\ell$, which may be zero. The bridge and the
typed-restriction convention are formalized in
\Cref{lem:bridge-decode}; the direct-query key-hit predicate ignores
the adversary's requested output length and is formalized as
$\KeyedTWOut$ in \Cref{lem:bridge-decode}.

The random oracle games are obtained from
\Cref{sec:games-pp,sec:mu-setting} by this mechanical substitution: direct
primitive access $(\pi, \pi^{-1})$ becomes finite-output access to $\ROflat$,
real construction or deterministic challenger checks use the corresponding
${}^\RO$ algorithm, and the ideal oracles $(\Rand,\Replay)$ are unchanged.
We use the same advantage names with subscript $\RO$: $\AdvColl_{\RO}$,
$\AdvPreSlice{a}{\mu}_{\RO}$, $\AdvePre_{\RO}$,
$\Advcmtds{\ell}_{\RO}$, $\Advcmts{\ell}_{\RO}$, and
$\AdvintRUP_{\RO}$ are the win probabilities of the corresponding random
oracle games, with CMTs challengers comparing deterministic encryptions via
$\Enc^\RO$. The all-in-one AE advantage $\Adv^{\mathsf{ae}}_{\RO}$ is the
corresponding real-vs-ideal distinguishing advantage with direct $\ROflat$
access in both worlds.

For mu-ind, $\Adv^{\mathsf{mu\text{-}ind}}_{\RO}(\mathcal{B})$ denotes the
\Cref{sec:mu-setting} game with the same $\mathsf{New}$ interface and per-user
nonce discipline: user $d$ uses
$(\Enc_{K_d}^{\RO}, \Ver_{K_d}^{\RO})$ when $b = 1$ and $(\Rand,\Replay)$ when
$b = 0$. Equivalently, in the two-game notation of \Cref{sec:mu-ind},
\[
  \Adv^{\mathsf{mu\text{-}ind}}_{\RO}(\mathcal{B}) =
  \bigl|
    \Pr[\mathcal{B}^{\mathsf{Real}^{\mathsf{mu}}_{\RO},\, \ROflat} \Rightarrow 1]
    - \Pr[\mathcal{B}^{\mathsf{Ideal}^{\mathsf{mu}},\, \ROflat} \Rightarrow 1]
  \bigr|,
\]
where the superscripts name the entire multi-user construction
interface in the corresponding world.
Adversary queries to $\ROflat$ are finite-output and counted in
$\qRO(\mathcal{B})$ of \Cref{sec:resources}. Construction-internal
$\RF(\cdot)$ lookups (those made by $\Enc_K^\RO$, $\Dec_K^\RO$,
$\PDec_K^\RO$, $\Ver_K^\RO$, or a hash or committing challenger) are not.

\subsection{Resources}
\label{sec:resources}

Adversary resources are measured by fixed, execution-uniform caps on
the number and size of oracle queries. An execution-uniform cap is one
that holds on every execution path, including over adversary
randomness, oracle randomness, and adaptive query choices.
Proofs may also introduce local execution-uniform caps, such as
$q_e$ for the number of accepted encryption queries, solely to index a
finite hybrid sequence. Such local caps are part of the adversary's
fixed query bounds, but they are not listed as headline resource parameters
unless they affect the final bound.

\Cref{tab:resource-accounting} gives the compact notation summary used in later
theorems. Lowercase resource quantities are realized counts in a random oracle
experiment, while uppercase quantities are execution-uniform caps used in the
public-permutation bounds. The detailed definitions below fix exactly which
computations contribute to each count, and \Cref{sec:loss-terms} defines the
closed-form loss terms.

\begin{table}[t]
  \centering
  \footnotesize
  \setlength{\tabcolsep}{3pt}
  \begin{tabular}{@{}p{0.18\linewidth}p{0.42\linewidth}p{0.32\linewidth}@{}}
    \toprule
    Symbol & Meaning & Primary use / definition \\
    \midrule
    $N$
      & Construction-internal \Keccakp{} calls induced by valid construction
        oracle queries and deterministic challenger checks.
      & Public-permutation lift and concrete work cap. \\
    $\Nprim$
      & Direct adversary queries to the public primitive interface.
      & Key-guess, root-collision, root-preimage terms; direct-query part of
        the lift. \\
    $\Ntot = N + \Nprim$
      & Converted primitive-query cost used by the flat sponge replacement.
      & Argument of $\Lift_c(\Ntot)$. \\
    $\qRO(\mathcal{B})$, $\QRO$
      & Direct finite-output queries to $\ROflat$ in the intermediate random
        oracle model; construction lookups are excluded.
      & Random oracle bounds; after the lift $\qRO \leq \Nprim$. \\
    $\qv(\mathcal{A})$, $\Qv$
      & Valid verification queries that are not replay-table hits.
      & Verification-tag terms in AE, mu-ind, and \textsf{INT-RUP} bounds. \\
    $\nleaf(\mathcal{B})$, $\Nleaf$
      & Distinct construction-generated final leaf transcripts.
      & Leaf tag collision and tag-guessing terms; always at most $N$. \\
    $D$
      & Execution-uniform cap on users activated in the mu-ind game.
      & Multi-user key-collision and key-guess terms. \\
    \midrule
    $\Lift_c(\Ntot)$
      & \cite{BDPV08} flat sponge differentiating loss.
      & Defined in \cref{sec:loss-terms}. \\
    $\KeyGuess_k(Q)$
      & Ideal-system key-guessing term.
      & Defined in \cref{sec:loss-terms}. \\
    $\RootColl_\tau(Q, s)$
      & $s$-way root-tag collision term.
      & Defined in \cref{sec:loss-terms}. \\
    $\RootPre_\tau(Q)$
      & Root-tag preimage term.
      & Defined in \cref{sec:loss-terms}. \\
    $\LeafColl_\tau(Q)$
      & Known-key leaf-tag collision term.
      & Defined in \cref{sec:loss-terms}. \\
    \bottomrule
  \end{tabular}
  \caption{Resource caps and recurring loss terms. Uppercase resource symbols
    are the caps that appear in public-permutation theorems; lowercase symbols
    are realized counts in random oracle executions or source games.}
  \label{tab:resource-accounting}
  \label{tab:notation}
\end{table}

\begin{description}
\item[$N$:] the construction interface cost, measured as the number of
  \Keccakp{} calls induced after the validity convention has been applied.
  It counts at per-duplex-call granularity, summing each root initialization,
  associated data duplex block, root message duplex block, leaf initialization,
  leaf message duplex block, and aggregation duplex block.

  \emph{AE-style games.} In real games, $N$ counts the valid construction
  computations actually run by the adversary's oracle queries. In ideal games,
  the same accounting applies to the corresponding real \TW{} computation on
  the valid query, even when the ideal oracle answers by table lookup or without
  running the construction. Valid plaintext-computing computations and valid
  AE-style non-replay verification computations are counted whether they accept
  or reject; replay-table verification hits are transcript-consistency checks
  and are not charged to $N$ or to the verification-tag terms.

  \emph{Hash games.} There is no adversary construction oracle. For the $\HTW$
  collision game, $N$ counts the $s$ deterministic encryptions on the submitted
  inputs. For the standard preimage game it counts both challenger encryptions:
  the source encryption computing $T^* = \HTW(X^*)$ and the deterministic
  encryption on the submitted input. For the everywhere-preimage game it counts
  the deterministic encryption on the submitted input. For the Sec and eSec
  games of \Cref{cor:second-preimage-rs04}, $N$ and $\Nprim$ are those of the
  explicit induced two-way collision adversary in that corollary, whose source
  and candidate deterministic encryptions are counted exactly as in the
  $s = 2$ collision game.

  \emph{Committing games.} For CMTDs-$\ell$, $N$ uses conservative
  committing-game accounting: after the syntax, domain, and message-length
  checks, it charges the deterministic decryption checks on the common
  ciphertext whether or not the particular $\WiC_\ell$ distinctness check would
  reject before those checks. Likewise, for the CMTs-$\ell$ source games of
  \Cref{cor:committing-combined}, after syntax and domain checks $N$ charges
  the deterministic encryptions that would be used to compare the $s$
  ciphertexts whether or not the $\WiC_\ell$ distinctness check would reject
  earlier.
\item[$\Nprim$:] the number of direct adversary queries to the
  primitive interface, the pair $(\pi, \pi^{-1})$ in
  public permutation games.
\item[$\Ntot = N + \Nprim$:] the total primitive-query cost
  used by the flat sponge lift. The lift of
  \Cref{app:flat-sponge-lift} is applied at $\Ntot$ because after the
  operational duplex-trace accounting of
  \Cref{lem:operational-duplex-trace}, the construction-internal duplex
  calls counted by $N$ and the adversary's direct primitive-interface
  queries counted by $\Nprim$ form a primitive-query sequence of cost at
  most $N+\Nprim$.
\item[$\qv(\mathcal{A})$, $\Qv$:] the number of verification queries
  with valid inputs that are not replay-table hits, and a cap on it.
  Used in the verification terms of the direct mu-ind, AE, and
  \textsf{INT-RUP} bounds of \Cref{sec:aead-security}. Every counted
  verification query contributes at least
  one verification/decryption
  construction call in the $N$ accounting, so $\qv(\mathcal{A}) \leq N$.
\item[$\qRO(\mathcal{B})$, $\QRO$:] the number of direct
  finite-output queries to $\ROflat$ by a random oracle model adversary
  $\mathcal{B}$, and a cap on it. Repeated direct queries on the same
  $Y$ are counted again. Construction-internal $\RF(\cdot)$ lookups are
  not counted. After the derived-adversary construction, $\qRO$ counts only
  simulator-induced $\ROflat$ calls from the adversary's direct
  primitive-interface queries, not the construction calls already
  counted in $N$ for the flat sponge lift.
\item[$\nleaf(\mathcal{B})$, $\Nleaf$:] the number of distinct
  construction-generated final leaf transcripts in the execution of
  $\mathcal{B}$, and a cap on it. Distinct final leaf transcripts are
  generated by valid encryption computations, by valid plaintext-computing
  computations, by valid AE-style
  verification computations (whether they accept or reject), and by
  hash and committing challenger construction computations that pass the
  early validity checks, including the preimage game's source encryption.
\end{description}

For a byte string $X$,
\[
  \leaves(X) = \max\bigl(0,\,
    \lceil \|X\| / \beta \rceil - 1\bigr)
\]
counts the leaf chunks induced by a body byte string of length $\|X\|$.
Thus $\nleaf$ is bounded by the sum of $\leaves(X)$ over the valid
construction computations that create final leaf transcripts, where
$X$ is the plaintext body for encryption computations and the submitted
ciphertext body for plaintext-computing, verification, and committing
decryption checks.
Each distinct final leaf transcript contains a final leaf duplex call
counted in $N$, so $\nleaf(\mathcal{B}) \leq N$ unconditionally. The
ratio $\nleaf/N$ is largest when most processed bytes lie in leaf
chunks, where each full leaf contributes $1 + \beta/\rho = 50$ duplex
calls to $N$ per final leaf transcript.

All theorems are stated for adversaries whose execution-uniform caps
hold at the stated bounds.

\subsection{Loss Terms}
\label{sec:loss-terms}

Five closed-form expressions recur across the security bounds: $\Lift_c$
throughout, with $\KeyGuess_k$, $\RootColl_\tau$, $\RootPre_\tau$, and
$\LeafColl_\tau$ in some. We define them here, after the resource caps that
parameterize them.

\paragraph{Flat sponge loss.}
The flat sponge differentiating loss is
\[
  \Lift_c(\Ntot) = \Ntot(\Ntot + 1) / 2^{c + 1}.
\]
For $\Ntot < 2^c$ it upper-bounds the \cite{BDPV08} sponge-differentiating
advantage; \Cref{lem:lift-bound} carries out the derivation from the
\cite[Lemmas~4 and~5]{BDPV08} product bounds through the Weierstrass
product inequality. For $\Ntot \geq 2^c$, $\Lift_c(\Ntot) \geq 1$ and the
advantage bound is trivial without any appeal to indifferentiability, and
already at $\Ntot \geq 2^{c/2}$ the bound exceeds $1/2$ and provides no
meaningful security.

\paragraph{Key-guessing term.}
The ideal system key-guessing term of \cite{BDPV11} is
\[
  \KeyGuess_k(Q) = Q / 2^k.
\]
It accounts for direct $\ROflat$ queries whose inputs satisfy the
decoder $\KeyedTWOut$ of \Cref{tab:keyed-twout}. Each such query is an
adaptive test of the decoded candidate key against the hidden key $K$,
independent of the query's requested output length. \Cref{lem:key-tests}
bounds the probability that any of $Q$ direct queries tests the right
key by $Q / 2^k$.

\paragraph{Root-collision term.}
The $s$-way root tag multi-collision term is
\[
  \RootColl_\tau(Q, s) = \binom{Q + s}{s} \big/ 2^{\tau (s - 1)},
\]
with the abbreviation $\RootColl_{\tauroot}(Q, s) =
\RootColl_\tau(Q, s)$ at $\tau = \tauroot$. The $+s$ is candidate-sequence
slack for the $s$ final root transcript inputs generated by the
challenger after the adversary outputs its tuples. These lookups are
construction-internal and are not added to $\qRO(\mathcal{B})$; after the
flat sponge lift of \Cref{app:flat-sponge-lift}, the public-permutation
bounds instantiate $Q$ as $\Nprim$.

\paragraph{Root-preimage term.}
The root tag preimage term is
\[
  \RootPre_\tau(Q) = (Q + 1) \big/ 2^{\tau},
\]
with the abbreviation $\RootPre_{\tauroot}(Q) = \RootPre_\tau(Q)$ at
$\tau = \tauroot$. It is the preimage analog of $\RootColl_\tau$: a
single query-independent target tag replaces the $s$-way internal
collision, one challenger lookup replaces the $s$ final-root lookups,
and the exponent drops from $\tau(s-1)$ to $\tau$. As in the
root-collision term, the flat sponge lift gives $Q = \Nprim$.

\paragraph{Leaf-collision term.}
The known-key leaf tag collision term is
\[
  \LeafColl_\tau(Q) = \binom{Q}{2} \big/ 2^\tau.
\]
It is the pairwise birthday bound on $\tau$-bit leaf tag projections
used by the collision-resistance theorem (\Cref{thm:coll}). In this game
the keys are adversary-supplied, so direct primitive queries can pre-read
candidate leaf tag points; after the flat sponge lift this gives
$Q=\Nprim+\Nleaf$. The hidden-key AE-style bounds instead carry the
linear same-slot term $\nleaf/2^{\tauleaf}$ of
\Cref{lem:leaf-collision}, written inline.

\paragraph{Tightness notes.}
The loss terms above are mostly generic search or birthday terms, and their
leading scales are tight in the usual random oracle sense. The key-guessing
term is matched at constants by direct random guesses for a $k$-bit key. In
the multi-user theorem's free-parameter setting, the term
$D\,\KeyGuess_k(\Nprim)$ is also tight at leading scale: using the per-user
nonce discipline of \Cref{sec:mu-ind}, which permits one nonce value across
users, an adversary obtains one known nonempty encryption from each targeted
user and compares each direct candidate-key query against all collected
root-initialization keystream blocks at once. Each guess then matches some
targeted user with probability $D/2^k$. The concrete joint cap charges the
target encryptions to $N$, lowering the feasible scale as discussed in
\Cref{sec:bound-discussion}.

The root terms are the standard random oracle search terms: an $s$-way
birthday search matches the leading scale of
$\RootColl_\tau(Q,s)$, and fixed-target preimage search matches the per-query
$2^{-\tau}$ rate of $\RootPre_\tau(Q)$. The hidden-key same-slot leaf term
$\nleaf/2^{\tauleaf}$ is also tight at constants for nonce-respecting
verification: after encrypting a two-chunk message under nonce $U$, an
adversary can submit non-replay verification queries that replace the leaf
ciphertext chunk by fresh strings of the same length while keeping the root
tag. Each query recomputes the hidden leaf tag at the same slot $(K,U,1)$,
and acceptance occurs when the fresh leaf tag equals the original hidden leaf
tag, with probability $2^{-\tauleaf}$ per verification-side final leaf
transcript.

The flat sponge lift term
$\Lift_c(\Ntot)=\Ntot(\Ntot+1)/2^{c+1}$ dominates the BDPV differentiating
bound $f_P(\Ntot)$ of \cite[Lemma~5]{BDPV08} via \Cref{lem:lift-bound}. That
differentiating bound is matched at leading order by the standard
inner-collision distinguisher; the remaining conservatism is in composing
indifferentiability into the keyed games, not in an obviously improvable
simulator bound. Real-vs-ideal games pay the same differentiating bound once
more, at $\Nprim$, to replace the ideal-side simulator by the real permutation
(\Cref{lem:sim-perm-proximity}); since $\Nprim \leq \Ntot$, that term never
dominates the $\Ntot$ term.
